\begin{exercise}{29}
A category $\mathcal C$ is called connected if
$\pi_0^\Delta(\mathrm N(\mathcal C))$ consists of only one element.
Show that a non-empty groupoid $\mathcal G$\footnote{The original
says ``a groupoid $\mathcal G$'', without the non-emptiness
hypothesis. Insert ``non-empty'': for the empty groupoid the
condition on every pair of objects is vacuously true, but
$\pi_0^\Delta(\mathrm N(\varnothing))=\varnothing$, which is not
a singleton.}
is connected if and only if for every two objects $x,y\in\mathcal G$,
the set $\operatorname{Hom}_{\mathcal G}(x,y)$ is non-empty.
Show that a connected groupoid is equivalent to $\mathbb B G$ for a group
$G$. Show, however, that the category of connected groupoids is
\emph{not} equivalent to the category of groups.

\begin{proof}
In a groupoid, the relation
$x\sim y\Longleftrightarrow\operatorname{Hom}_{\mathcal G}(x,y)\ne\varnothing$
is already an equivalence relation, by identities, inverses and
composition. Thus Definition~1.1.9 identifies
$\pi_0^\Delta(\mathrm N(\mathcal G))$ with its set of classes.
For a non-empty groupoid this is a singleton precisely when
every pair of objects is related.

Now suppose $\mathcal G$ is connected. Choose an object $x$ and put
$G=\operatorname{Aut}_{\mathcal G}(x)$.
The inclusion $I\colon\mathbb B G\to\mathcal G$ sends its sole
object to $x$. Choose arrows $\gamma_y\colon x\to y$ with
$\gamma_x=\mathrm{id}_x$. A quasi-inverse sends every object to
the sole object of $\mathbb B G$ and sends $f\colon y\to z$ to
\[
 R(f)=\gamma_z^{-1}f\gamma_y.
\]
Cancellation of $\gamma_z\gamma_z^{-1}$ shows that $R$ is a
functor. We have $RI=\mathrm{id}$, and the arrows $\gamma_y$
form a natural isomorphism $IR\Rightarrow\mathrm{id}_{\mathcal G}$.
Hence $\mathcal G\simeq\mathbb B G$.

Finally, the category of groups has an initial object, the
trivial group. The category of connected groupoids and functors
has none: let $D$ have objects $0,1$ and exactly one morphism
between each ordered pair. It is a connected groupoid, and any
connected groupoid $\mathcal H$ admits two distinct functors
$\mathcal H\to D$, constant at $0$ and at $1$ respectively.
They are distinct because $\mathcal H$ is non-empty.
Thus no $\mathcal H$ can be initial. Since equivalences preserve
initial objects, these two categories are not equivalent.
\end{proof}
\end{exercise}
